\documentclass[11pt]{article}
\usepackage[T1]{fontenc}
\usepackage{lmodern,amsmath,amssymb,mathtools,microtype,booktabs}
\usepackage[letterpaper,top=.92in,bottom=.9in,left=1.12in,right=1.12in]{geometry}
\usepackage{xcolor,fancyhdr}
\usepackage[colorlinks=true,linkcolor=teal!65!black,citecolor=teal!60!black,urlcolor=teal!60!black,bookmarksnumbered=true]{hyperref}
\hypersetup{pdftitle={Sharp moment bounds for the first zero and the zeros encoded in their errors},pdfauthor={OpenAI}}
\definecolor{ink}{HTML}{183E4D}
\pagestyle{fancy}\fancyhf{}\lhead{\small\color{ink}Sharp moment bounds for the first zero}\rhead{\small\color{ink}Research note}\cfoot{\small\thepage}
\setlength{\headheight}{15pt}
\setlength{\emergencystretch}{2em}
\newcommand{\repo}[2]{\href{#1}{#2}}
\begin{document}
\thispagestyle{plain}
{\small\sffamily\color{ink} ANALYTIC NUMBER THEORY \hfill 26 NOVEMBER 2026}\par\vspace{12pt}
{\huge\bfseries\color{ink}Sharp moment bounds for the first zero\par}
\vspace{3pt}{\LARGE\color{ink}and the zeros encoded in their errors\par}
\vspace{13pt}
\noindent\textbf{Abstract.}
From $d$ high-order power sums of inverse squared zero distances we construct a
polynomial whose distinguished root bounds the nearest distance. For fixed
$d$ the side alternates with parity. The leading error contains the product
of the next $d$ normalized inverse squared distances, with an explicit
factor for ties. Differences of successive bounds therefore recover higher
zero distances exponentially in the stride, including in the presence of
multiple zeros. When the first $d-1$ tail entries are separated, the bound
is the exact endpoint consistent with those $d$ moments and the prescribed
leading multiplicity. We construct the unique extremizer and exhibit a
repeated-tail example where endpoint attainment fails. Dirichlet GRH supplies
the real-zero input for the arithmetic applications.
\medskip

\noindent\textbf{A first case.}
Write $0<\gamma_1<\gamma_2\le\cdots$ for the positive zeta ordinates
with a simple first zero and put $p_k=\sum_j\gamma_j^{-2k}$. Once
$p_k^2<2p_{2k}$, the two elementary estimators are
\[
      R_{1,k}=p_k,\qquad
      R_{2,k}=\frac{p_k+\sqrt{2p_{2k}-p_k^2}}{2}.
\]
For large $k$, $R_{2,k}\le\gamma_1^{-2k}\le R_{1,k}$. If the second
ordinate occurs $m$ times, then
$R_{1,k}-R_{2,k}\sim m\gamma_2^{-2k}$. The analogous discrepancy at
the next order involves both $\gamma_2$ and $\gamma_3$. The two-moment
extremum already belongs to the classical theory of sharp norm inequalities
for probability vectors~\cite{sakai}. We treat arbitrary orders, classify a
sharp case and quantify the errors.

\paragraph{Connection with the supplied sources.}
The public declarations in Sernau's Lean project give Mathlib's RH and the
critical-line theorem for every Dirichlet character in the open strip.
Its guide traces a one-sided prime-remainder estimate to three inputs about
short forms, spanning and a determinant lower bound. We use the resulting
zero location at the point where signs of symmetric sums become decisive.
The supplied encyclopedia article records the formal-proof history and
relates zero location to prime-number errors. The Lean FRO and Mathlib
Initiative certified the original RH statement and proof, and Sutherland
published independent verification records~\cite{lean,verification,wiki}.
Our arguments below are ordinary mathematical proofs. We inspected the
upstream source and its verification record without rebuilding it.

\section{Reading the spectrum from alternating errors}
\label{sec:deflation}

Let $a=x_1=\cdots=x_r>x_{r+1}\ge x_{r+2}\ge\cdots>0$,
with $\sum_i x_i<\infty$. Entries are listed with their multiplicities.
The list can be infinite or finite. Write $p_k=\sum_i x_i^k$ for $k\ge1$.
Every quantity used to define the estimators below depends only on these power sums.
For each $k$ put
\[
    h_k=\left\lfloor\frac{p_k^2}{p_{2k}}\right\rfloor,
    \qquad c_k=\frac{p_k}{h_k}.
\]
The integer $h_k$ estimates the multiplicity $r$. For $d\ge1$, form the degree-$d$
polynomial
\begin{equation}\label{eq:defpoly}
 P_{d,k}(t)=[u^d]\frac{\prod_i(1+x_i^k u)}{(1+tu)^{h_k}} .
\end{equation}
Here $[u^d]$ extracts a Taylor coefficient at zero. To evaluate
\eqref{eq:defpoly} from power sums, use Newton's identities
\[
 e_{0,k}=1,\qquad
 j e_{j,k}=\sum_{\ell=1}^j(-1)^{\ell-1}p_{k\ell}e_{j-\ell,k},\qquad
 P_{d,k}(t)=\sum_{j=0}^d(-1)^j{h_k+j-1\choose j}e_{d-j,k}\,t^j.
\]
Choose the real root of $P_{d,k}$ nearest $c_k$ when this choice exists and
is unique, and call it $R_{d,k}$. The theorem ensures that every fixed order
$d$ eventually has such a root.

\paragraph{Theorem 1 (alternating errors with multiplicities).}
Suppose the tail $x_{r+1},x_{r+2},\ldots$ has at least $d$ entries. Write
$q_j=x_{r+j}/a$, $b_d=\prod_{j=1}^d q_j$, and let $C_d$ be the number of
$d$-element subsets of tail indices with product $b_d$ in the normalized
tail. Explicitly, if $v$ tail entries strictly exceed $q_d$ and $w$ equal
$q_d$ (using normalized values), then $C_d={w\choose d-v}$.
For all sufficiently large $k$, $h_k=r$ and $R_{d,k}$ exists uniquely
under the selection rule above. It satisfies
\begin{equation}\label{eq:defasym}
 \frac{R_{d,k}}{a^k}-1
    =(-1)^{d+1}\frac{C_d}{r}\,b_d^k\bigl(1+o(1)\bigr).
\end{equation}
Thus an even order gives an eventual lower bound for $a^k$ and an odd order an
eventual upper bound. The factor controlling the error contains one factor for
each of the first $d$ tail entries, including repetitions.

\paragraph{Proof.}
Scaling every $x_i$ by $1/a$ leaves $h_k$ unchanged and divides each root and
$c_k$ by $a^k$, by homogeneity. We may therefore take $a=1$. Put
$B_k=\sum_{j\ge1}q_j^k$ and $D_k=\sum_{j\ge1}q_j^{2k}$. Both tend to zero
by dominated convergence. Also $0\le D_k\le B_k$, so that
\[
       (r+B_k)^2-r(r+D_k)=r(2B_k-D_k)+B_k^2\ge0.
\]
Consequently $p_k^2/p_{2k}\to r$ from above (with equality for an empty tail),
and its floor is eventually $r$. In particular $c_k\to1$.

For fixed $d$, the coefficient of $u^j$ in the tail product
$\prod_i(1+q_i^k u)$ is at most $B_k^j/j!$ for $j\ge1$. Therefore
$P_{d,k}(t)$ converges coefficientwise, as do its derivatives, to
\[
 H_d(t)=[u^d]\frac{(1+u)^r}{(1+tu)^r},\qquad
 H_d(1)=0,\qquad H_d'(1)=-r[u^{d-1}](1+u)^{-1}=(-1)^d r.
\]
On a sufficiently small fixed interval $[1-\epsilon,1+\epsilon]$, the derivative
of $H_d$ has constant strict sign and the endpoint values of $H_d$ have
opposite signs. Uniform convergence on this interval gives the same properties
for $P_{d,k}$ for large $k$. There is exactly one root there, and shrinking the
interval shows that this root tends to $1$. All other real roots stay outside
the original interval, so $c_k\to1$ identifies the indicated root uniquely.

At the true value $t=1$ the dominant factor cancels exactly:
\[
 P_{d,k}(1)=
 \sum_{i_1<\cdots<i_d}(q_{i_1}\cdots q_{i_d})^k
          = C_d b_d^k\bigl(1+o(1)\bigr).
\]
Indeed, after division by $b_d^k$, each maximizing subset contributes one and
every other subset tends to zero. For $k\ge1$ the normalized terms are
dominated by $q_{i_1}\cdots q_{i_d}/b_d$, whose sum is finite. The
description of $C_d$ follows by taking every entry strictly larger than the
cutoff, then choosing the remaining entries at the cutoff.
By the real mean value theorem,
$R_{d,k}-1=-P_{d,k}(1)/P'_{d,k}(\xi_k)$ for a point $\xi_k$ between
$1$ and $R_{d,k}$. Its denominator tends to $(-1)^d r$.
Rescaling gives \eqref{eq:defasym}. \hfill$\square$

\paragraph{Corollary 2 (recovery from successive differences).}
For an infinite list define
\[
 L_0=\lim_{k\to\infty}p_k^{1/k},\qquad
 L_d=\lim_{k\to\infty}|R_{d,k}-R_{d+1,k}|^{1/k}\quad(d\ge1).
\]
Then
\begin{equation}\label{eq:defrecover}
   L_0=a,\qquad L_d=a b_d,\qquad x_{r+d}=L_0\frac{L_d}{L_{d-1}}.
\end{equation}
The formula holds for finite lists too whenever the tail contains at least
$d$ entries. If there are exactly $d$, the order-$d+1$ polynomial vanishes
at $a^k$ identically once $h_k=r$, and its nearest root is eventually $a^k$.

\paragraph{Proof.}
Since $0<q_{d+1}<1$, the relative error at order $d+1$ is $o(b_d^k)$.
Subtraction gives
$|R_{d,k}-R_{d+1,k}|=(C_d/r)(a b_d)^k(1+o(1))$.
For a finite tail of length $d$, the same calculation uses exact cancellation
at order $d+1$ and the simple limiting root proved above. Taking roots and
ratios gives \eqref{eq:defrecover}, using $p_k/a^k\to r$.
\hfill$\square$

The input at each fixed $(d,k)$ consists of finitely many moments. The direction
of the bound and the multiplicity choice are eventual statements. When $k$ is
small, a value of $h_k$ alone need not certify a multiplicity. Numerically,
successive orders can differ by much less than machine precision.


\subsection*{Exponential recovery from two consecutive strides}
For an infinite list put $G_{0,k}=p_k$ and
$G_{d,k}=|R_{d,k}-R_{d+1,k}|$ for $d\ge1$. Define
$D_{d,k}=G_{d,k+1}/G_{d,k}$ once the selection rules are unique and the
denominator is positive. Then
\begin{equation}\label{eq:fast}
 \widehat{x}_{r+d,k}=D_{0,k}\frac{D_{d,k}}{D_{d-1,k}}
       = x_{r+d}\bigl(1+O(\eta_d^k)\bigr)
       \quad\text{for some }0<\eta_d<1.
\end{equation}
This removes the factor $C_d/r$ even when several zeros have the same distance.
No earlier recovered entry is needed as input.

To verify the rate, let $\rho_j$ be the ratio of the next distinct $j$-fold
tail product to the largest, with value zero if all products are maximal.
For fixed $j$, only finitely many products exceed any positive threshold
because they have finite sum. Thus $\rho_j<1$. Set $\rho_0=0$.
The total normalized contribution of nonmaximal products to $P_{j,k}(1)$
in the proof of Theorem~1 is $O(\rho_j^{k-1})$. If $\rho_j=0$, this
contribution is zero. Moreover $\sum_i q_i^k=O(q_1^k)$, so all coefficients
of $P_{j,k}$ and its derivative differ from their respective limits by
$O(q_1^k)$ in the normalized units. The point in the mean value theorem
lies within $O(b_j^k)$ of one. The derivative there is consequently
$(-1)^j r+O(q_1^k)$. Subtracting the order-$j+1$ approximation introduces
an additional relative error $O(q_{j+1}^k)$. We obtain
\[
 G_{j,k}=\frac{C_j}{r}(a b_j)^k
      \{1+O(\max(q_1,\rho_j)^k)\}.
\]
Taking consecutive-stride ratios proves \eqref{eq:fast} with
$\eta_d=\max(q_1,\rho_d,\rho_{d-1})$. In particular, near coincidences of
distinct distances can slow convergence even when the closest zero is
well separated.

There is a simpler way to recover ordered entries if no bound on the first
one is wanted. For fixed $j$, the same dominant-product argument gives
$e_{j,k}^{1/k}\to x_1\cdots x_j$. The point of \eqref{eq:fast} is that
its input consists of discrepancies between one-sided bounds for the
\emph{same} leading entry. The next section identifies conditions under which
each bound is optimal for the moment data at its degree.

\section{Which bound is best possible from the given moments?}

Fix $d\ge2$ and suppose the input consists of $r$ and exactly the numbers
$p_k,p_{2k},\ldots,p_{dk}$. Here a competitor is any finite or countable multiset
of positive summable numbers with these same moments and with largest entry of
multiplicity $r$. Its cardinality may vary. We optimize the $k$th power of its
largest entry. The restriction on $r$ matters when $r>1$. If
$p_k^2<2p_{2k}$, every competitor already has $r=1$.

\paragraph{Theorem 3 (sharp attainable endpoint).}
Assume the original tail has at least $d$ entries and
\[
     q_1>q_2>\cdots>q_{d-1}>q_d .
\]
For sufficiently large $k$, the largest possible powered entry among the
competitors is $R_{d,k}$ when $d$ is odd. The smallest possible powered entry
is $R_{d,k}$ when $d$ is even. Exactly one multiset attains the stated
endpoint. It has $r$ equal leading entries and $d-1$ distinct smaller entries.
If admissible multisets must be infinite, the same value is the supremum or
infimum, respectively.

\paragraph{Proof.}
Scale to $a=1$ and take $k$ large enough that $h_k=r$. If $T$ is a competing
largest powered entry, then
\begin{equation}\label{eq:momentbracket}
             p_{2k}/p_k\ \le\ T\ \le\ p_k/r .
\end{equation}
Indeed $p_{2k}\le T p_k$ and the $r$ leading entries contribute $rT$ to $p_k$.
Both endpoints tend to one. The polynomial $P_{d,k}$ is common to every
competitor because Newton's identities use the stated data. At $T$ its value
is the degree-$d$ elementary symmetric sum of the remaining powered entries.
In particular it is nonnegative. By Theorem~1, on a fixed neighborhood of one
containing both $R_{d,k}$ and the entire interval \eqref{eq:momentbracket},
$(-1)^dP'_{d,k}>0$ for large $k$ (in the proof, this follows from the simple
limiting zero). Hence $T\le R_{d,k}$ for odd $d$ and $T\ge R_{d,k}$ for even $d$.

To attain equality, put $R=R_{d,k}$, $m=d-1$ and
\[
 B_j(t)=[u^j]\frac{\prod_i(1+x_i^ku)}{(1+tu)^r},
 \qquad b_j^*=B_j(R),\qquad H(u)=\sum_{j=0}^{m}b_j^*u^j.
\]
Since $B_d(R)=0$, $(1+Ru)^rH(u)$ agrees with the original product through
degree $d$. It remains to show that $H(u)$ factors as
$\prod_{i=1}^m(1+z_i u)$ with $0<z_i<R$.

Write $b_0=1$ as well as $b_j=q_1\cdots q_j$ for $j\ge1$. For
$1\le j\le m$, separation makes the maximal $j$-fold tail product unique,
so $B_j(1)\sim b_j^k$. For fixed $j$, $B'_j$ is bounded uniformly on a
neighborhood of one by coefficient convergence. Theorem~1 gives
$|R-1|=O(b_d^k)$, and $b_d<b_j$. It follows that $b_j^*\sim b_j^k$.
To see where the factors lie, consider
\[
      V(z)=\sum_{j=0}^m(-1)^j b_j^* z^{m-j}.
\]
At $z=c q_i^k$ with fixed $c>0$, the terms with indices $i-1$ and $i$
dominate. Indeed the ratio between the absolute values of consecutive terms
at indices $j-1,j$ is asymptotic to $(q_j/q_i)^k/c$. After dividing $V(z)$
by $(-1)^{i-1}b_{i-1}^*z^{m-i+1}$ its value tends to $1-1/c$.
Taking $c=1-\epsilon$ and $1+\epsilon$, with $0<\epsilon<1$, supplies a
positive root in each of $m$ eventually disjoint intervals. Thus $V$ has $m$
distinct roots $z_i\sim q_i^k$, all eventually smaller than $R$.
This proves the desired factorization. The powered multiset consisting of
$r$ copies of $R$ and the $z_i$ has the required moments by Newton's
identities. Take positive $k$th roots to return to the original units.

At equality the degree-$d$ symmetric sum of the competitor's tail is zero.
A positive tail with $d$ entries would give a positive sum. The tail
therefore has at most $d-1$ entries, and its full product must equal $H$
by the same coefficient identities. This proves uniqueness up to ordering.

For infinite competitors perturb the $d$ distinct powered locations of this
finite extremizer. The Jacobian of their first $d$ power sums has entries
$\ell m_j z_j^{\ell-1}$, where the locations now include $R$ and $m_j$ are
their multiplicities. Its determinant is nonzero by the Vandermonde formula.
Append the powered entries $\epsilon 2^{-n}$ for $n\ge1$. Their vector of
first $d$ power sums tends to zero as $\epsilon\downarrow0$. The inverse
function theorem permits adjustments of the original $d$ locations to restore
those moments exactly. Positivity, strict order and top multiplicity persist,
and the added $k$th roots have finite sum. Letting $\epsilon$ tend to zero
proves the last assertion. \hfill$\square$

\paragraph{A finite test.}
The preceding proof also yields a certificate at a specified $k$. With $r$
known, or certified as one, construct $P$ using $r$ in its denominator.
Locate a real root $R$ and an interval containing both that root and
\eqref{eq:momentbracket} on which $(-1)^dP'>0$. This suffices for the claimed
one-sided bound. Certify that the residual $V$ above has $d-1$ distinct roots
in $(0,R)$ to prove that the bound is attained. Both tasks use only polynomial
signs. The script \texttt{check\_endpoint.py} illustrates these checks using
rational arithmetic. Separation is a sufficient condition for their eventual
success. Theorem~1 itself permits arbitrary ties throughout the tail.

\paragraph{The role of equal weights.}
Each occurrence contributes one copy of its powered value to the moments.
An arbitrary positive measure with freely variable weights defines a different
extremal problem. In the zeta application the integer counts come from zero
multiplicities. The synthetic spectra realizing Theorem~3 need not be spectra
of any Dirichlet $L$-function. The sharpness concerns what follows from these
moments and their multiplicities alone.

\paragraph{A genuine obstruction at a tie.}
For the spectrum $(1,1,q,q,q)$, with $0<q<1$, take $d=3$ and put $w=q^k$.
Theorem~1 gives $R_{3,k}=1+O(w^3)$. At that root the first two coefficients
of $\prod_i(1+x_i^ku)/(1+R_{3,k}u)^2$ equal $3w+O(w^3)$ and
$3w^2+O(w^3)$, respectively. Indeed the values at $t=1$ are exactly
$3w,3w^2$, and their derivatives in $t$ are bounded near $1$. If a
competitor attained this upper endpoint, the residual sum $S$ and sum of
squares $Q$ after deletion of its two largest entries would satisfy
$S^2/Q\to3$ as $k\to\infty$. But for every residual list of nonnegative
entries $y_i$, Cauchy--Schwarz gives
$\sum y_i^3\ge Q^2/S$, whence
\[
  e_3(y)\ge \frac{(S^2-Q)(S^2-2Q)}{6S}.
\]
For each sufficiently large fixed $k$ this lower bound is strictly positive
at the proposed endpoint, whereas $P_{3,k}(R_{3,k})=0$. It stays positive in
a neighborhood by continuity, while $P_{3,k}(t)$ tends to zero there as
$t\to R_{3,k}$. Thus even the supremum over competitors with leading
multiplicity two is strictly below the order-three upper bound in this
example. The asymptotic error law still holds with its prescribed tie factor.

\section{Local Taylor data of completed $L$-functions}

Let $F(z)$ be an entire function of order at most one, all of whose zeros
$\gamma$ are real. Fix $\tau\in\mathbb R$ with $F(\tau)\ne0$ and list
$x_i=|\gamma_i-\tau|^{-2}$ in decreasing order, counting multiplicities.
Hadamard factorization gives
\begin{equation}\label{eq:defhad}
    G_{\tau}(v)=
    \frac{F(\tau+\sqrt v)F(\tau-\sqrt v)}{F(\tau)^2}
          =\prod_i(1-x_i v),\qquad
    p_n=-n[v^n]\log G_{\tau}(v).
\end{equation}
Indeed, apply the genus-one factorization to $F(\tau+z)$. Its factors
$(1-z/(\gamma-\tau))\exp(z/(\gamma-\tau))$ pair at $z,-z$ to give
$1-z^2/(\gamma-\tau)^2$. The exponential prefactors cancel after
normalization, and $\sum_i|\gamma_i-\tau|^{-2}<\infty$ by the order bound.
The numerator is even in $\sqrt v$, so this construction is independent
of the chosen square root. The logarithm is the local branch with value zero
at $v=0$.

For the Riemann zeta function take
$F(z)=\xi(\tfrac12+iz)$, where
$\xi(s)=\tfrac12s(s-1)\pi^{-s/2}\Gamma(s/2)\zeta(s)$.
For a primitive nonprincipal character $\chi$ of conductor $Q$, let
$\kappa\in\{0,1\}$ satisfy $\chi(-1)=(-1)^\kappa$ and take
\[
 F(z)=\Lambda(\tfrac12+iz,\chi),\qquad
 \Lambda(s,\chi)=
       (Q/\pi)^{(s+\kappa)/2}\Gamma((s+\kappa)/2)L(s,\chi).
\]
These completions are entire of order one. Under the respective critical-line
theorem their zeros are real in $z$, so \eqref{eq:defhad} supplies every
moment required by Theorem~1. This argument allows a complex character and an
arbitrary real observation point away from a zero. The recovered multiplicity at
distance $D$ counts all zeros at $\tau-D$ and $\tau+D$ together. Thus a tie
can arise from two distinct ordinates. In the zeta case at $\tau=0$, one may
instead use the even function $\xi(\tfrac12+i\sqrt v)/\xi(\tfrac12)$.
It is the normalized square root of $G_0(v)$ and counts each positive ordinate
once for each occurrence, without also counting its reflected partner.

For the positive zeta ordinates there is also an especially convenient
Taylor coordinate at $s=1$. The even product for $\xi$ gives
\[
 \frac{\xi(1+u)}{\xi(1)}
 =\prod_{\gamma>0}\left(1+\frac{u+u^2}{1/4+\gamma^2}\right).
\]
Indeed $(1/2+u)^2-1/4=u+u^2$. Thus the coefficients in $v=u+u^2$ of
$\log(\xi(1+u)/\xi(1))$ yield the power sums of $(1/4+\gamma^2)^{-1}$
after multiplication in degree $j$ by $(-1)^{j-1}j$.

\subsection*{A finite certificate from four moments}
The order-three and order-four signs also give an exact finite example.
Use the coordinate $y=200/(1/4+\gamma^2)$ for the positive zeta ordinates
and let $M_j=\sum y^j$, with multiplicity. The ProveIt packet
\texttt{li\_ritz\_zero\_recovery.zip} computes enclosures for these moments
from the Taylor series of $\log\xi$ at $1$, using an Euler--Maclaurin
remainder and outward-rounded arithmetic. We reran its moment generator
and reproduced its data byte for byte. With stride four, the only inputs
used here are $M_4,M_8,M_{12},M_{16}$. In that order they lie between
$n/10^{26}$ and $(n+1)/10^{26}$ for the following integers:
\[
\begin{array}{rr}
105572466326784184430480813,&100025018857591281580900012,\\
99765036478443525027857289,&99677225165408692087482695.
\end{array}
\]
Put $S=M_4$, $Q=M_8$, and let $A=\max y^4$. The inequalities
$Q/S\le A\le S$ place $A$ in $I=[0.94,1.18]$.
The strict inequality $S^2/Q<2$ proves the nearest positive ordinate
is simple. For $d=3,4$ construct $P_d=P_{d,4}$ with $h=1$.
Direct interval evaluation certifies $P'_3<0$ and $P'_4>0$ throughout
$I$. Let $t(g)=(200/(1/4+g^2))^4$.
Both $t(14.13472107)$ and $t(14.13472515)$ belong to $I$, and exact
rational arithmetic gives
\[
   P_3(t(14.13472107))<0,\qquad P_4(t(14.13472515))<0.
\]
Since $P_d(A)\ge0$ for both orders, monotonicity of $P_d$ and $t$ gives
\[
             \boxed{14.13472107<\gamma_1<14.13472515.}
\]
The included \texttt{zeta\_finite.py} checks the two signs and tests the
derivatives on sixteen equal subintervals of $I$ by interval Horner
evaluation. The certificate uses RH to identify all the Taylor moments
with positive real ordinate data. Its purpose is to demonstrate the
sparse-moment bounds. The ordinate and its simplicity are classical,
and the earlier packet provides tighter bounds using other moments.

At this stride one can certify optimality as well. Each $P_d$ has a simple
root $T_d$ in $I$. The checker proves
\[
\begin{aligned}
0.999193608630<T_3<0.999193608633,\\
0.999191309425<T_4<0.999191309428.
\end{aligned}
\]
The interval $I$ contains every real number whose distance to $S$ is at
most $|S-T_d|$, as follows directly from the input intervals. Thus the
uniqueness on $I$ also verifies the nearest-root selection rule.
Write $B_j=[u^j]\prod_i(1+y_i^4u)/(1+T_du)$.
Then $B_0=1$, $B_d=0$ and $B_j=e_{j,4}-T_dB_{j-1}$.
For $d=3$, alternating endpoint signs put the two zeros of
$\sum_{j=0}^{d-1}(-1)^j B_j z^{d-1-j}$ in $(16,17)/1000$ and
$(39,40)/1000$. For $d=4$ the three intervals are $(5,6)/1000$,
$(8,9)/1000$, and $(41,42)/1000$. All these possible tail entries are
positive and smaller than $T_d$. Appending them to $T_d$ gives a list
with the prescribed first $d$ power sums, by Newton's identities.
Every competing list has its maximum in $I$ and satisfies $P_d(A)\ge0$.
Thus $T_3$ is the largest possible powered maximum given the first three
stride moments, and $T_4$ is the smallest given the first four. These
extremal values range over positive spectra sharing the stated data.

\subsection*{Numerical illustration for the Riemann zeros}
For a separate experiment in the coordinate $x=200/\gamma^2$, we evaluated
$\xi(1/2+\sqrt{200z})/\xi(1/2)=\prod_{\gamma>0}(1+200 z/\gamma^2)$
on a circle, extracted its coefficients by discrete Cauchy quadrature, and
used Newton's identities with $x_j=200/\gamma_j^2$. The scale $200$ keeps the
largest entry near one. Ordinates enter the program only after the estimates,
in a separate reference calculation.

\begin{center}
\begin{tabular}{r r r r r}\toprule
Stride & First, order 3 & Second, gaps & Third, gaps & Fourth, gaps\\\midrule
4  &14.134721124&20.216005&24.320741&28.744971\\
8  &14.134725142&20.938539&25.011406&29.975006\\
16 &14.134725142&21.019510&25.012261&30.353284\\
24 &14.134725142&21.021935&25.010938&30.410957\\
32 &14.134725142&21.022035&25.010862&30.421911\\\midrule
Reference&14.134725142&21.022040&25.010858&30.424876\\\bottomrule
\end{tabular}
\end{center}

The last three columns use Corollary~2, replacing $L_0$ by
$(p_k/h_k)^{1/k}$ and $L_j$ by $G_{j,k}^{1/k}$. At stride 32, the
order-three estimate reproduces the first ordinate to the 17 significant
digits printed by the script. The fourth ordinate is recovered with absolute
error about $0.0030$. The recorded output also gives the consecutive-stride
variant \eqref{eq:fast}. It cancels the tie constant, though its error at
moderate strides need not be smaller. We obtained the displayed numbers at
150 decimal digits with 192 nodes, and repeated the computation at 185 digits
with 240 nodes. These are floating-point experiments. The next example
isolates polynomial roots using exact rational arithmetic.

\section{An exact-arithmetic example with repeated entries}
\label{sec:example}

Take the finite spectrum
\[
  (1,1,\tfrac12,\tfrac12,\tfrac12,\tfrac13,\tfrac14).
\]
The leading multiplicity is $r=2$, and the first tail value has multiplicity
three. Hence $(b_1,b_2,b_3,b_4)=(1/2,1/4,1/8,1/24)$ and
$(C_1,C_2,C_3,C_4)=(3,3,1,1)$. The table gives
$(-1)^{d+1}r(R_{d,k}-1)/(C_d b_d^k)$, whose predicted limit is one in
each column.
\begin{center}
\begin{tabular}{r|rrrr}
$k$ & $d=1$ & $d=2$ & $d=3$ & $d=4$ \\ \hline
 8 & 1.014308231 & 1.055560592 & 1.142759979 & 1.125152057 \\
16 & 1.000512566 & 1.001583578 & 1.004659175 & 1.010114634 \\
32 & 1.000000773 & 1.000002319 & 1.000006955 & 1.000100454
\end{tabular}
\end{center}
The included \texttt{check\_rational.py} computes the moments and polynomial
coefficients as rational numbers. It isolates each selected root on an interval
centered at $c_k$, proving that the derivative has fixed sign there and checking
opposite endpoint signs. Exact bisection then encloses every normalized value
in an interval of width below $10^{-18}$.


\section{Comparison and further questions}
The ProveIt index describes 61 deposited packets as filed without independent
checks at placement. Its nested incoming directory contains a closer
predecessor, \emph{Recovering individual zeta zeros from Li coefficients}.
That paper uses $x=(1/4+\gamma^2)^{-1}$ and a growing Hankel matrix pencil
to bound individual zeros. Its Ritz estimates approach each fixed spectral
point as matrix dimension increases. Here the polynomial degree is fixed,
the moment stride grows and all the polynomial roots we select estimate the
same first spectral point. Their errors encode the rest. Another indexed
packet studies separation of number fields through differences of Li
coefficients~\cite{proveit}. The accompanying source inventory identifies
the scope of the comparison.

Voros developed the secondary zeta functions of zero ordinates.
Kawalec expresses successive ordinates by limits that subtract the already
known zeros~\cite{voros,kawalec}. Graeffe iteration amplifies separation
by powering polynomial roots~\cite{graeffe}, and leading elementary symmetric
coefficients already detect successive products. Anari et al. study
worst-case approximation of a largest root from leading coefficients;
Acevedo et al. describe the boundary geometry of power-sum images~\cite{anari,acevedo}.
The contribution proposed here is the multiplicity-dependent alternating
error formula together with the exact endpoint, its attaining spectrum, and
the recovery rule from differences of consecutive endpoints when the separation
hypotheses hold. The comparison supports a distinction from the inspected
results; historical priority remains open.

\paragraph{Further research questions.}
For tied tail entries the inequality in Theorem~1 remains valid while its
extremal interpretation can fail. Which multiplicity patterns determine the
sharp endpoint in this case? How rapidly may $d$ grow with $k$ while the
root-selection and error formulas stay uniform? Given rigorous absolute
errors in the input moments, what choice of stride and degree gives the
narrowest certified zero interval? The cancellation in successive differences
makes the last question relevant even when the exact moments uniquely
determine the required limit.

\begin{thebibliography}{10}\small
\bibitem{lean} L. Sernau et al., \emph{GRH for Dirichlet $L$-functions, in Lean 4}.
\repo{https://github.com/lsernau/grh-lean/tree/72fead102f663d5c7a407964a13dc0cf9d977207}{Repository at 72fead1},
especially \texttt{grh/src/Statements.lean}, \texttt{grh/src/Main.lean}
and \texttt{grh/GUIDE.md}.
\bibitem{verification} Lean FRO and Mathlib Initiative,
\repo{https://leanprover-community.github.io/blog/posts/riemann-hypothesis/}{\emph{Lean verification of the Riemann Hypothesis}}, 21 November 2026.
A. Sutherland, \repo{https://math.mit.edu/~drew/grh/}{verification files and independent kernel checks}.
\bibitem{wiki} Wikipedia contributors, \repo{https://en.wikipedia.org/wiki/Riemann_hypothesis}{\emph{Riemann hypothesis}}, accessed 26 November 2026.
\bibitem{proveit} V. Reshetnikov, \emph{ProveIt}, \repo{https://github.com/VladimirReshetnikov/ProveIt/tree/244c77760d4dd675372c5bda2438c135e2fc802d/Analysis/GeneralizedRiemannHypothesis}{Generalized Riemann Hypothesis collection, snapshot 244c777}.
See \repo{https://github.com/VladimirReshetnikov/ProveIt/blob/244c77760d4dd675372c5bda2438c135e2fc802d/Analysis/GeneralizedRiemannHypothesis/docs/incoming/li_ritz_zero_recovery.zip}{\texttt{li\_ritz\_zero\_recovery.zip}}
and \texttt{Research/grh\_method\_transfers\_research}.
\bibitem{voros} A. Voros, \repo{https://doi.org/10.5802/aif.1955}{Zeta functions for the Riemann zeros}, \emph{Ann. Inst. Fourier} 53 (2003), 665--699.
\bibitem{kawalec} A. Kawalec, \repo{https://arxiv.org/abs/2012.06581}{Analytical recurrence formulas for non-trivial zeros of the Riemann zeta function}, arXiv:2012.06581.
\bibitem{graeffe} G. Malajovich and J. P. Zubelli, \repo{https://arxiv.org/abs/math/9908149}{On the geometry of Graeffe iteration}, \emph{J. Complexity} 17 (2001), 541--573.
\bibitem{anari} N. Anari, S. Oveis Gharan, A. Saberi and N. Srivastava,
\repo{https://arxiv.org/abs/1704.03892}{Approximating the largest root and applications to interlacing families}, SODA 2018.
\bibitem{acevedo} J. Acevedo, G. Blekherman, S. Debus and C. Riener,
\repo{https://arxiv.org/abs/2303.09512}{The wonderful geometry of the Vandermonde map},
\emph{Found. Comput. Math.} 26 (2026), 2005--2051.
\bibitem{sakai} Y. Sakai and K. Iwata, \repo{https://arxiv.org/abs/1605.00019}{Sharp bounds between two R\'enyi entropies of distinct positive orders}, arXiv:1605.00019.
\end{thebibliography}
\end{document}
